\usepackage{xeCJK}
\usepackage{geometry}
\geometry{left=2cm,right=2cm,top=2.5cm,bottom=2.5cm}
\setlength{\headheight}{12.7pt}

\usepackage{titlesec}
\titleformat{\section}
  {\normalfont\large\bfseries}
  {\thesection}
  {1em}
  {}
\titlespacing*{\section}
  {0pt}
  {3.5ex plus 1ex minus .2ex}
  {2.3ex plus .2ex}

\usepackage{amsmath}
\usepackage{amssymb}
\usepackage{amsthm}
\usepackage{enumitem}
\setlist[enumerate]{itemsep=0pt,topsep=0pt,parsep=0pt,leftmargin=0pt,itemindent=2.5em}
\setlist[itemize]{itemsep=0pt,topsep=0pt,parsep=0pt,leftmargin=0pt,itemindent=2.5em}
\usepackage{fancyhdr}
\pagestyle{fancy}
\fancyhf{}
\usepackage{graphicx}
\usepackage{hyperref}
\usepackage{indentfirst}
\usepackage{lmodern}


\begin{document}
\setlength{\abovedisplayskip}{1pt}
\setlength{\belowdisplayskip}{1pt}

\section{序数算术与字典序}

序数算术恰好是简单运算的\href{../01 序关系/05 字典序#nameddest=sec:定义1.1}{字典序}.

\vspace{20pt}

\hypertarget{sec:定理1.1}{}
\noindent\textbf{定理1.1}

任给\href{../02 序数/01 序数#nameddest=sec:定义1.1}{序数} $\alpha,\beta$, 有
$\alpha+\beta=\mathrm{Type}(\alpha\sqcup\beta,\leqslant_{\mathrm{lex}})$.

\noindent{\kaishu 证明.}

对任意的$(i,x)\in \alpha\sqcup\beta$, 定义
\[
    f(i,x)=\alpha i+x,
\]
易证$f:\alpha\sqcup\beta\to\alpha+\beta$是双射, 下证$f$严格保序.

在$\alpha\sqcup\beta$中任取$(i,x)<_{\mathrm{lex}}(j,y)$ :
\begin{enumerate}
    \item 如果$i=0,\,j=1$, 则$f(i,x)=x<\alpha\leqslant\alpha+y=f(j,y)$ ;
    \item 如果$i=j=0,\,x<y$, 则$f(i,x)=x<y=f(j,y)$ ;
    \item 如果$i=j=1,\,x<y$, 则$f(i,x)=\alpha+x<\alpha+y=f(j,y)$.
    \hfill $\qedsymbol$
\end{enumerate}

\vspace{20pt}

\hypertarget{sec:定理1.2}{}
\noindent\textbf{定理1.2}

任给序数$\alpha,\beta$, 有
$\alpha\beta=\mathrm{Type}(\beta\times\alpha,\leqslant_{\mathrm{lex}})$.

\noindent{\kaishu 证明.}

对任意的$(x,y)\in\beta\times\alpha$, 定义
\[
    f(x,y)=\alpha x+y,
\]
易证$f:\beta\times\alpha\to\alpha\beta$是双射, 下证$f$严格保序.

在$\beta\times\alpha$中任取$(x,y)<_{\mathrm{lex}}(x',y')$ :
\begin{enumerate}
    \item 如果$x<x'$, 则
    \[
        f(x,y)=\alpha x+y<\alpha x+\alpha=\alpha x^+
        \leqslant\alpha x'\leqslant\alpha x'+y'=f(x',y')\,;
    \]
    \item 如果$x=x',\,y<y'$, 则
    \begin{equation*}
        f(x,y)=\alpha x+y<\alpha x'+y'=f(x',y').
    \tag*{\qedsymbol}\end{equation*}
\end{enumerate}

\vspace{30pt}

\hypertarget{sec:定理1.3}{}
\noindent\textbf{定理1.3}

任给序数$\alpha,\beta$ :
\begin{enumerate}
    \item $\alpha+\beta\approx\alpha\sqcup\beta\approx
    \beta\sqcup\alpha\approx\beta+\alpha$,
    \item $\alpha\beta\approx\beta\times\alpha\approx
    \alpha\times\beta\approx\beta\alpha$.
\end{enumerate}

\end{document}